\documentclass[10pt,a4paper]{amsart}
\usepackage[margin=19mm]{geometry}
\usepackage{amsmath,amssymb,mathtools}
\usepackage[scr=rsfs,cal=euler]{mathalfa}
\usepackage{xfrac}
\usepackage[unicode,hidelinks,bookmarks=false]{hyperref}
\newcommand{\ie}{i.e.\ }
\newcommand{\cf}{cf.\ }
\newcommand{\resp}{respectively\ }
\newcommand{\Subsection}{\subsection}
\providecommand{\nobreakdash}{\nobreak\hbox{-}}
\setlength{\emergencystretch}{2em}
\multlinegap0pt
\usepackage{enumerate}
\usepackage{amssymb}
\usepackage{xcolor}
\newtheorem{proposition}{Proposition}
\DeclareMathOperator{\PSU}{PSU}
\DeclareMathOperator{\Soc}{Soc}
\title{A Determination of $B$-groups of Order $p^4$}
\author{Christopher Herbig}
\date{Source: arXiv:2608.26463v1 (CC BY 4.0)}
\begin{document}
\maketitle

\noindent\emph{Source and licence.} A Determination of $B$-groups of Order $p^4$. Version arXiv:2608.26463v1 (CC BY 4.0), distributed by arXiv under the Creative Commons Attribution 4.0 International licence, http://creativecommons.org/licenses/by/4.0/, as recorded in the arXiv licence metadata for this exact version and shown on the abstract page https://arxiv.org/abs/2608.26463v1. Full licence notice: \texttt{licenses/arxiv-2608.26463-CC-BY-4.0.txt}.\par\smallskip

\noindent\textit{Editorial scope.} Counted unit: the article's proposition proposition (source0.tex lines 528--530). The article's complete original proof of it is delivered with it (source0.tex lines 531--549, one \texttt{proof} environment of 19 lines). No mathematics is written, repaired or re-derived: the statement and the proof are the article's own lines, in the article's own notation, and no retained line is substituted or re-flowed.  No further source-local context is needed: every cross-reference of every retained line resolves inside the two delivered spans.  Nothing was omitted from any retained span, so the package carries no omission record.  No retained line contains a citation, so the wrapper carries no bibliography and no bibliography entry.  No retained line contains a figure environment, an image-inclusion command or drawing code, so no image is delivered and the package declares no asset dependency.  Editorial formatting only, in addition: an AMS \texttt{amsart} preamble that restates the article's own declarations and macros used by the retained lines, namely the environments \texttt{proposition} on separate counters, together with the standard packages those lines need.  The retained environments print in this document as Proposition 1 in order of appearance, whereas the article prints the same items with the numbering of its own full document.  The complete native source of the article as distributed in the arXiv e-print for this exact version is bundled unchanged as \texttt{sources/208635/source0.tex}.
\begin{proposition}
	Let $G$ be an almost simple group which acts faithfully and uniprimitively on a set $\Omega$ where $|\Omega| = p^a$ for some prime $p$ and an integer $a$. If $\alpha \in \Omega$, then $|G:G_\alpha| = 27$, and $\Soc(G) \cong \PSU(4,2)$.
\end{proposition}	

\begin{proof}
	Let $G$ be an almost simple group, and let $G$ act faithfully and primitively on a set $\Omega$ where $|\Omega| = p^a$. Also, set $S = \Soc(G)$. Clearly, we may assume that $G$ is nonsimple, for otherwise, the statement reduces to the above result of Guralnick. Choose $\alpha \in \Omega$. Since $G$ acts faithfully, $S$ is not a subgroup of $G_\alpha$. Conversely, $G_\alpha$ cannot be a proper subgroup of $S$, since $G_\alpha$ is necessarily maximal by primitivity. Consequently, $SG_\alpha = G$. By the second isomorphism theorem, we have that \newline $|S:S\cap G_\alpha| = p^a$. 
	
	By Guralnick's result above, the action of $S$ on the cosets of $S \cap G_\alpha$ is doubly transitive unless $S \cong \PSU(4,2)$ and $p^a = 27$. Under the latter circumstance, there is nothing to prove, so we can assume that $S$ acts doubly transitively on the cosets of $S \cap G_\alpha$. It now suffices to show that the action of $G$ on the cosets of $G_\alpha/G$ is doubly transitive.
	
	Consider the map $(S \cap G_\alpha)s \longmapsto G_\alpha s$ from $(S \cap G_\alpha)/S$ to $G_\alpha/G$. It is straightforward to verify that this map is bijective and does not depend on a choice of coset representatives. In particular, a set of coset representatives for $S$ over $S \cap G_\alpha$ is a set of coset representatives for $G$ over $G_\alpha$, and for $s_1, s_2 \in S$, we have $G_\alpha s_1 = G_\alpha s_2$ iff $(S \cap G_\alpha)s_1 = (S \cap G_\alpha)s_2$.  Now, given two pairs of cosets, say $(G_\alpha s_1, G_\alpha s_2)$ and $(G_\alpha s_3, G_\alpha s_4)$ for $s_i \in S$ and $i=1, 2, 3, 4$, we want to find $g \in G$ such that 
	$$
		(G_\alpha s_1 g, G_\alpha s_2 g) = (G_\alpha s_3, G_\alpha s_4).
	$$
	By the double transitivity of $S$, we can find $s \in S$ such that
	$$
		((S \cap G_\alpha) s_1 s, (S \cap G_\alpha) s_2 s) = ((S \cap G_\alpha) s_3, (S \cap G_\alpha) s_4).
	$$
	From this, we obtain
	$$
		(G_\alpha s_1 s, G_\alpha s_2 s) = (G_\alpha s_3, G_\alpha s_4),
	$$
	and it follows that $G$ acts doubly transitively on the cosets of $G_\alpha$ and therefore also acts doubly transitively on $\Omega$.
\end{proof}

\end{document}
